\section{Related work and position of this paper}\label{sec:related}

Four literatures meet in the problem we study: detecting manipulation, finding coordination among traders, testing coincidence between event series, and using simulated markets and enforcement theory to choose rules. For each we say what it provides, what it assumes, and what is left for this paper; Table~\ref{tab:positioning} summarises.

\paragraph{Detection with labelled events or order-book data.} Most detection work targets a manipulation that leaves a public trace or a label. For pump-and-dump schemes organised on messaging channels the announcement itself labels the event: \citet{xu2019} study 412 pumps announced on Telegram (June 2018 to February 2019) and predict which coin will be pumped; \citet{hu2022} model 709 events from 2019 to 2022 with a sequence network over channel histories; \citet{lamorgia2020} propose a real-time detector and \citet{lamorgia2023} reach an F1 score of 94.5\% with detection 25 seconds after a pump starts, on a data set of about 900 events; \citet{dhawan2023} measure average price distortions of 65\% on 355 cases; \citet{kamps2018} define criteria and apply anomaly detection to trade data. For spoofing, \citet{tuccella2021} train a recurrent network on granular order-book data for early detection, and \citet{fabre2025} predict the distribution of mid-price moves from Hawkes-process order-flow variables and find that 31\% of the large orders in a three-day sample of Level-3 data could spoof. Wash trading is detected from structure: loops of matched orders in directed trader graphs on stock data \citep{cao2016}, accounts and trading structures that meet legal definitions on two blockchain exchanges \citep{victor2021}, and statistical regularities of reported volume, which put wash trading above 70\% of reported volume on unregulated exchanges \citep{cong2023}. These methods need labels (announcements, cases) or counterparty-level records, and the spoofing and pump settings are not those of the Vietnamese cases, in which groups of accounts trade among themselves or continuously for months (Appendix~\ref{app:cases}). We use a pump data set only to measure how weak the preparation phase of a pump is in public tick data (Section~\ref{sec:pump}); our detector uses neither labels nor counterparties, only the timing and side of submissions.

\paragraph{Coordination among traders.} Closer to our setting, collusion has been sought in trade networks. \citet{palshikar2008} cluster the trade graph and combine evidence from several algorithms; \citet{sun2012} test the correlation between the degree and strength of traders in networks of eight manipulated and 44 non-manipulated stocks against randomised networks; \citet{jaeger2025} link insiders by the temporal similarity of their reported trades in 2.9 million filings, motivated precisely by the lack of labels. \citet{saavedra2011} show that synchrony is a common trait of honest traders: day traders who trade more in step with others lose less, so synchrony alone does not signal collusion. That finding is the reason for our emphasis on honest groups (market-making desks, order-splitting funds, bot fleets) and for the netting of buys against sells (Proposition~\ref{prop:netting}), and it is borne out on real order flow (Section~\ref{sec:real-accounts}). Relative to this strand we compare one account with the \emph{aggregate} of all others rather than pairs, which pools evidence when a ring fragments its trading over many accounts, and we attach an exact null to the comparison.

\paragraph{Coincidence tests with surrogates.} Testing whether events in two series coincide more than chance without a model for either series is a long-standing problem in neuroscience. \citet{grun2009} show that nonstationary rates and non-Poisson structure generate false positives unless they are part of the null; \citet{louis2010} build surrogates that keep rates and inter-event intervals; \citet{amarasingham2012} review jitter resampling and exact tests for repeated fine-scale patterns; \citet{platkiewicz2017} prove that interval jitter gives an exact test while spike-centred jitter can inflate false positives; \citet{liang2025} give finite-sample false-positive guarantees for time-shift tests that allow dependence within the series. We transfer the circular-shift version to accounts, add netting, take the maximum over tolerances inside one shift distribution, and prove two properties that matter for a regulator: the resolution limit of the discrete $p$-value across many accounts (Proposition~\ref{prop:budget}) and a power law in the number of accounts (Proposition~\ref{prop:power}). The warning of \citet{grun2009} is the main lesson of our real-data test: common drivers make the null false for most active accounts.

\paragraph{Simulated markets for rules and for surveillance.} Open simulators exist \citep{byrd2020}, and their realism is a recognised issue: \citet{vyetrenko2019} compare simulated and real limit order books on stylised facts and find more realistic behaviour when the fundamental comes from historical data, and \citet{platt2018} find that microstructure parameters of agent-based models can be calibrated but behavioural parameters remain problematic. Both point to the design we adopt, a family of markets matching the observable moments instead of a single tuned market (Section~\ref{sec:br}). Agent-based models have been used to study rules: tick size, which worsens quality, especially for small caps, as it grows \citep{zhao2020}; price limits, which cause volatility spillover and delay price discovery when upper limits bind \citep{zhang2016}; and minimum resting times, circuit breakers, cancellation fees and taxes, which trade stability against resilience in a market with high-frequency traders \citep{leal2019}. For manipulation specifically, \citet{wang2020} generate order streams in which a manipulator mimics a market maker, and \citet{wang2018} analyse a cloaking rule against a spoofing exploiter by empirical game-theoretic analysis; \citet{cartea2020} show in a control model that a larger fine makes spoofing less attractive until it disappears. These studies concern single manipulators and generic markets. We study a ring, a specific venue, and the best response of the manipulator to each rule in a family of calibrated markets, add closed-form conditions (deterrence capacity, block-price law), and verify the mechanism ex post.

\paragraph{Enforcement economics and Vietnamese evidence.} Deterrence depends on the probability of detection and the size of the sanction \citep{becker1968,polinsky2000}; inspection games formalise limited inspection capacity \citep{avenhaus2002}, and double-oracle methods compute equilibria of games with an adversary \citep{mcmahan2003,lanctot2017}, a natural next step for the interaction between an adaptive ring and an adaptive regulator. \citet{aggarwal2006} use SEC cases to show that manipulation raises volatility, liquidity and returns during the episode, with prices falling afterwards; \citet{comerton2014} estimate that about 1\% of closing prices are manipulated, that only a small fraction is detected and prosecuted, and that the regulator's budget has a strong effect on both manipulation and detection, which is the empirical counterpart of our capacity result (Proposition~\ref{prop:deter}); \citet{khwaja2005} study price manipulation by intermediaries in an emerging market. For Vietnam, \citet{nguyen2017} study insider trading around stock splits, \citet{nguyen2023} herding across COVID-19 periods, and \citet{vo2023} find that the 2016 tick reduction lowered trading costs on the Ho Chi Minh Stock Exchange, a result we use as an external check (Section~\ref{sec:tick}). None of these studies has account-level manipulation data for Vietnam, and neither do we.

\begin{table}[t]
\centering\footnotesize
\begin{adjustbox}{max width=\linewidth}
\begin{tabular}{p{3.6cm}p{3.4cm}p{1.4cm}p{1.6cm}p{2.3cm}p{1.8cm}}
\toprule
strand & typical input & labels needed & account-level data & error control & policy evaluation\\
\midrule
pump-and-dump detection & public tick data and announcements & yes & no & empirical & no\\
spoofing detection & order book, orders & yes & no (one trader) & empirical & little\\
wash-trading detection & counterparty-level trades & rules or labels & yes & rules, empirical & no\\
trader-network collusion & trade graphs, filings & some & yes & randomised networks & no\\
time-shift and jitter tests & event series (neuroscience) & no & not applicable & exact or finite-sample & no\\
agent-based rule studies & simulated markets & no & simulated & not applicable & yes, fixed agents\\
\textbf{this paper} & order submissions (time, side) & no & yes (not public here) & exact per account; ranking in practice & yes, best response over calibrated markets\\
\bottomrule
\end{tabular}
\end{adjustbox}
\caption{Positioning of the paper among the strands of related work.}\label{tab:positioning}
\end{table}

\paragraph{Position.} The gap this paper addresses is the combination of the strands in the last row: a screening statistic that needs neither labels nor counterparties, whose error properties and power are stated and checked, together with a policy evaluation in which the manipulator adapts and the market is not a single draw. Our claims are limited accordingly. Components are not new (shift tests, adversarial simulation, enforcement economics); the contribution is their calibrated combination, the theory that connects detection power to deterrence capacity and to the effect of a block-price rule, and the measurement of where the combination fails, on simulated honest groups and on real account-level order flow.
